% Part: intuitionistic-logic
% Chapter: tableaux
% Section: proofs

\documentclass[../../../include/open-logic-section]{subfiles}

\begin{document}

\olfileid{int}{tab}{prf}

\olsection{అంతఃప్రజ్ఞావాద తర్కానికి \usetoken{టాబ్లోలు}{tableau}}

\begin{ex}
  $(!A \land !B) \lif !C \Proves
  !A \lif (!B \lif !C)$ అని చూపే
  సంవృత టాబ్లోను ఇస్తున్నాం.
  \begin{oltableau}
    [\pFmla{\True}{(\formula{A} \land \formula{B}) \lif \formula{C}}{1},
      just =\TAss
      [\pFmla{\False}{\formula{A} \lif (\formula{B} \lif \formula{C})}{1},
        just =\TAss
        [\pFmla{\True}{\formula{A}}{1.1},
          just = {\TRule{\False}{\lif}[2]}
          [\pFmla{\False}{\formula{B} \lif \formula{C}}{1.1},
            just = {\TRule{\False}{\lif}[2]}
            [\pFmla{\True}{\formula{B}}{1.1.1},
              just = {\TRule{\False}{\lif}[4]}
              [\pFmla{\False}{\formula{C}}{1.1.1},
                just = {\TRule{\False}{\lif}[4]}
                [\pFmla{\False}{\formula{A} \land \formula{B}}{1.1.1},
                  just = {\TRule{\True}{\lif}[1]}
                  [\pFmla{\False}{\formula{A}}{1.1.1},
                    just= {\TRule{\False}{\land}[7]}, close]
                  [\pFmla{\False}{\formula{B}}{1.1.1},
                    just= {\TRule{\False}{\land}[7]}, close]]
                [\pFmla{\True}{\formula{C}}{1.1.1},
                  just = {\TRule{\True}{\lif}[1]}, close]]
            ]
          ]
        ]
      ]
    ]
  \end{oltableau}
  \sourcecorrection{OLTEINTTABPRF-001}{మూల టాబ్లోలో అసత్య సంయోగ విసర్జనతో వచ్చే రెండు శాఖలకు కారణంగా నాలుగో పంక్తిని చూపింది; సంయోగం ఉన్న ఏడో పంక్తికే ఆ నియమం వర్తిస్తుంది కాబట్టి రెండు సూచనలను ఏడో పంక్తికి సరిచేశాం.}
\end{ex}



\begin{prob}
  కింది వాటిని చూపే సంవృత అంతఃప్రజ్ఞావాద
  \tetoken{టాబ్లోలను}{!!{tableau}s} కనుగొనండి:
  \begin{enumerate}
    \item $\Proves !A \lif (!B \lif !A)$
    \item $\Proves \lnot(!A \land \lnot !A)$
    \item $!A \lif (!B \lif !C) \Proves (!A \land !B) \lif !C$
    \item $\lnot !A \lor \lnot !B \Proves \lnot(!A \land !B)$
  \end{enumerate}
\end{prob}

\end{document}
